
\documentclass{article}

\usepackage{a4wide}
\usepackage[utf8]{inputenc}
\usepackage{enumerate}
\usepackage{xspace}
\usepackage{longtable}
\usepackage{graphicx}
\usepackage{listings}
\usepackage{xcolor}
\usepackage{calc}
\usepackage{lmodern}
\usepackage{amssymb}
\usepackage{amsmath}
\usepackage{mathdots}
\usepackage{mathtools}
\usepackage{multicol}
\usepackage{enumitem}
\usepackage{tasks}
\usepackage{pdfpages}
\usepackage{listings}
\usepackage{bera}

\lstset{basicstyle=\ttfamily,
escapeinside={||},
mathescape=true}

\usepackage{ifthen}
\usepackage{amssymb}
\usepackage{multicol}
\usepackage{graphicx}
\usepackage[absolute]{textpos}
\usepackage{amsmath, amscd, amssymb, amsthm, latexsym}
% \usepackage[noload]{qtree}
%\usepackage{xspace,rotating,calligra,dsfont,ifthen}
\usepackage{xspace,rotating,dsfont,ifthen}
\usepackage[spanish,activeacute]{babel}
\usepackage[utf8]{inputenc}
\usepackage{pgfpages}
\usepackage{pgf,pgfarrows,pgfnodes,pgfautomata,pgfheaps,xspace,dsfont}
\usepackage{listings}
\usepackage{multicol}
\usepackage{todonotes}
\usepackage{url}
\usepackage{float}
\usepackage{framed,mdframed}
\usepackage{cancel}

\usepackage[strict]{changepage}


\makeatletter


\newcommand\hfrac[2]{\genfrac{}{}{0pt}{}{#1}{#2}} %\hfrac{}{} es un \frac sin la linea del medio

\newcommand\Wider[2][3em]{% \Wider[3em]{} reduce los m\'argenes
\makebox[\linewidth][c]{%
  \begin{minipage}{\dimexpr\textwidth+#1\relax}
  \raggedright#2
  \end{minipage}%
  }%
}


\@ifclassloaded{beamer}{%
  \newcommand{\tocarEspacios}{%
    \addtolength{\leftskip}{4em}%
    \addtolength{\parindent}{-3em}%
  }%
}
{%
  \usepackage[top=1cm,bottom=2cm,left=1cm,right=1cm]{geometry}%
  \usepackage{color}%
  \newcommand{\tocarEspacios}{%
    \addtolength{\leftskip}{3em}%
    \setlength{\parindent}{0em}%
  }%
}

\newcommand{\encabezadoDeProc}[4]{%
  % Ponemos la palabrita problema en tt
%  \noindent%
  {\normalfont\bfseries\ttfamily proc}%
  % Ponemos el nombre del problema
  \ %
  {\normalfont\ttfamily #2}%
  \
  % Ponemos los parametros
  (#3)%
  \ifthenelse{\equal{#4}{}}{}{%
  \ =\ %
  % Ponemos el nombre del resultado
  {\normalfont\ttfamily #1}%
  % Por ultimo, va el tipo del resultado
  \ : #4}
}

\newcommand{\encabezadoDeTipo}[2]{%
  % Ponemos la palabrita tipo en tt
  {\normalfont\bfseries\ttfamily tipo}%
  % Ponemos el nombre del tipo
  \ %
  {\normalfont\ttfamily #2}%
  \ifthenelse{\equal{#1}{}}{}{$\langle$#1$\rangle$}
}

% Primero definiciones de cosas al estilo title, author, date

\def\materia#1{\gdef\@materia{#1}}
\def\@materia{No especifi\'o la materia}
\def\lamateria{\@materia}

\def\cuatrimestre#1{\gdef\@cuatrimestre{#1}}
\def\@cuatrimestre{No especifi\'o el cuatrimestre}
\def\elcuatrimestre{\@cuatrimestre}

\def\anio#1{\gdef\@anio{#1}}
\def\@anio{No especifi\'o el anio}
\def\elanio{\@anio}

\def\fecha#1{\gdef\@fecha{#1}}
\def\@fecha{\today}
\def\lafecha{\@fecha}

\def\nombre#1{\gdef\@nombre{#1}}
\def\@nombre{No especific'o el nombre}
\def\elnombre{\@nombre}

\def\practicas#1{\gdef\@practica{#1}}
\def\@practica{No especifi\'o el n\'umero de pr\'actica}
\def\lapractica{\@practica}


% Esta macro convierte el numero de cuatrimestre a palabras
\newcommand{\cuatrimestreLindo}{
  \ifthenelse{\equal{\elcuatrimestre}{1}}
  {Primer cuatrimestre}
  {\ifthenelse{\equal{\elcuatrimestre}{2}}
  {Segundo cuatrimestre}
  {Verano}}
}


\newcommand{\depto}{{UBA -- Facultad de Ciencias Exactas y Naturales --
      Departamento de Computaci\'on}}

\newcommand{\titulopractica}{
  \centerline{\depto}
  \vspace{1ex}
  \centerline{{\Large\lamateria}}
  \vspace{0.5ex}
  \centerline{\cuatrimestreLindo de \elanio}
  \vspace{2ex}
  \centerline{{\huge Pr\'actica \lapractica -- \elnombre}}
  \vspace{5ex}
  \arreglarincisos
  \newcounter{ejercicio}
  \newenvironment{ejercicio}{\stepcounter{ejercicio}\textbf{Ejercicio
      \theejercicio}%
    \renewcommand\@currentlabel{\theejercicio}%
  }{\vspace{0.2cm}}
}


\newcommand{\titulotp}{
  \centerline{\depto}
  \vspace{1ex}
  \centerline{{\Large\lamateria}}
  \vspace{0.5ex}
  \centerline{\cuatrimestreLindo de \elanio}
  \vspace{0.5ex}
  \centerline{\lafecha}
  \vspace{2ex}
  \centerline{{\huge\elnombre}}
  \vspace{5ex}
}


%practicas
\newcommand{\practica}[2]{%
    \title{Pr\'actica #1 \\ #2}
    \author{Algoritmos y Estructuras de Datos I}
    \date{Segundo Cuatrimestre 2019}

    \maketitlepractica{#1}{#2}
}

\newcommand \maketitlepractica[2] {%
\begin{center}
\begin{tabular}{r cr}
 \begin{tabular}{c}
{\large\bf\textsf{\ Algoritmos y Estructuras de Datos I\ }}\\
Primer Cuatrimestre 2020\\
\title{\normalsize Gu\'ia Pr\'actica #1 \\ \textbf{#2}}\\
\@title
\end{tabular} &
\begin{tabular}{@{} p{1.6cm} @{}}
\includegraphics[width=1.6cm]{logodpt.jpg}
\end{tabular} &
\begin{tabular}{l @{}}
 \emph{Departamento de Computaci\'on} \\
 \emph{Facultad de Ciencias Exactas y Naturales} \\
 \emph{Universidad de Buenos Aires} \\
\end{tabular}
\end{tabular}
\end{center}

\bigskip
}


% Símbolos varios

\newcommand{\nat}{\ensuremath{\mathds{N}}}
\newcommand{\ent}{\ensuremath{\mathds{Z}}}
\newcommand{\float}{\ensuremath{\mathds{R}}}
\newcommand{\bool}{\ensuremath{\mathsf{Bool}}}
\newcommand{\True}{\ensuremath{\mathrm{true}}}
\newcommand{\False}{\ensuremath{\mathrm{false}}}
\newcommand{\Then}{\ensuremath{\rightarrow}}
\newcommand{\Iff}{\ensuremath{\leftrightarrow}}
\newcommand{\implica}{\ensuremath{\longrightarrow}}
\newcommand{\IfThenElse}[3]{\ensuremath{\mathsf{if}\ #1\ \mathsf{then}\ #2\ \mathsf{else}\ #3\ \mathsf{fi}}}
\newcommand{\In}{\textsf{in }}
\newcommand{\Out}{\textsf{out }}
\newcommand{\Inout}{\textsf{inout }}
\newcommand{\yLuego}{\land _L}
\newcommand{\oLuego}{\lor _L}
\newcommand{\implicaLuego}{\implica _L}
\newcommand{\cuantificador}[5]{%
	\ensuremath{(#2 #3: #4)\ (%
		\ifthenelse{\equal{#1}{unalinea}}{
			#5
		}{
			$ % exiting math mode
			\begin{adjustwidth}{+2em}{}
			$#5$%
			\end{adjustwidth}%
			$ % entering math mode
		}
	)}
}

\newcommand{\existe}[4][]{%
	\cuantificador{#1}{\exists}{#2}{#3}{#4}
}
\newcommand{\paraTodo}[4][]{%
	\cuantificador{#1}{\forall}{#2}{#3}{#4}
}

% Símbolo para marcar los ejercicios importantes (estrellita)
\newcommand\importante{\raisebox{0.5pt}{\ensuremath{\bigstar}}}


\newcommand{\rango}[2]{[#1\twodots#2]}
\newcommand{\comp}[2]{[\,#1\,|\,#2\,]}

\newcommand{\rangoac}[2]{(#1\twodots#2]}
\newcommand{\rangoca}[2]{[#1\twodots#2)}
\newcommand{\rangoaa}[2]{(#1\twodots#2)}

%ejercicios
\newtheorem{exercise}{Ejercicio}
\newenvironment{ejercicio}[1][]{\begin{exercise}#1\rm}{\end{exercise} \vspace{0.2cm}}
\newenvironment{items}{\begin{enumerate}[a)]}{\end{enumerate}}
\newenvironment{subitems}{\begin{enumerate}[i)]}{\end{enumerate}}
\newcommand{\sugerencia}[1]{\noindent \textbf{Sugerencia:} #1}

\lstnewenvironment{code}{
    \lstset{% general command to set parameter(s)
        language=C++, basicstyle=\small\ttfamily, keywordstyle=\slshape,
        emph=[1]{tipo,usa}, emphstyle={[1]\sffamily\bfseries},
        morekeywords={tint,forn,forsn},
        basewidth={0.47em,0.40em},
        columns=fixed, fontadjust, resetmargins, xrightmargin=5pt, xleftmargin=15pt,
        flexiblecolumns=false, tabsize=2, breaklines, breakatwhitespace=false, extendedchars=true,
        numbers=left, numberstyle=\tiny, stepnumber=1, numbersep=9pt,
        frame=l, framesep=3pt,
    }
   \csname lst@SetFirstLabel\endcsname}
  {\csname lst@SaveFirstLabel\endcsname}


%tipos basicos
\newcommand{\rea}{\ensuremath{\mathsf{Float}}}
\newcommand{\cha}{\ensuremath{\mathsf{Char}}}
\newcommand{\str}{\ensuremath{\mathsf{String}}}

\newcommand{\mcd}{\mathrm{mcd}}
\newcommand{\prm}[1]{\ensuremath{\mathsf{prm}(#1)}}
\newcommand{\sgd}[1]{\ensuremath{\mathsf{sgd}(#1)}}

\newcommand{\tuple}[2]{\ensuremath{#1 \times #2}}

%listas
\newcommand{\TLista}[1]{\ensuremath{seq \langle #1\rangle}}
\newcommand{\lvacia}{\ensuremath{[\ ]}}
\newcommand{\lv}{\ensuremath{[\ ]}}
\newcommand{\longitud}[1]{\ensuremath{|#1|}}
\newcommand{\cons}[1]{\ensuremath{\mathsf{addFirst}}(#1)}
\newcommand{\indice}[1]{\ensuremath{\mathsf{indice}}(#1)}
\newcommand{\conc}[1]{\ensuremath{\mathsf{concat}}(#1)}
\newcommand{\cab}[1]{\ensuremath{\mathsf{head}}(#1)}
\newcommand{\cola}[1]{\ensuremath{\mathsf{tail}}(#1)}
\newcommand{\sub}[1]{\ensuremath{\mathsf{subseq}}(#1)}
\newcommand{\en}[1]{\ensuremath{\mathsf{en}}(#1)}
\newcommand{\cuenta}[2]{\mathsf{cuenta}\ensuremath{(#1, #2)}}
\newcommand{\suma}[1]{\mathsf{suma}(#1)}
\newcommand{\twodots}{\ensuremath{\mathrm{..}}}
\newcommand{\masmas}{\ensuremath{++}}
\newcommand{\matriz}[1]{\TLista{\TLista{#1}}}

% Acumulador
\newcommand{\acum}[1]{\ensuremath{\mathsf{acum}}(#1)}
\newcommand{\acumselec}[3]{\ensuremath{\mathrm{acum}(#1 |  #2, #3)}}

% \selector{variable}{dominio}
\newcommand{\selector}[2]{#1~\ensuremath{\leftarrow}~#2}
\newcommand{\selec}{\ensuremath{\leftarrow}}

\newcommand{\pred}[3]{%
    {\normalfont\bfseries\ttfamily\noindent pred }%
    {\normalfont\ttfamily #1}%
    \ifthenelse{\equal{#2}{}}{}{\ (#2) }%
    \{%
    \begin{adjustwidth}{+2em}{}
      \ensuremath{#3}
    \end{adjustwidth}
    \}%
    {\normalfont\bfseries\,\par}%
}

\newenvironment{proc}[4][res]{%

  % El parametro 1 (opcional) es el nombre del resultado
  % El parametro 2 es el nombre del problema
  % El parametro 3 son los parametros
  % El parametro 4 es el tipo del resultado
  % Preambulo del ambiente problema
  % Tenemos que definir los comandos requiere, asegura, modifica y aux
  \newcommand{\pre}[2][]{%
    {\normalfont\bfseries\ttfamily Pre}%
    \ifthenelse{\equal{##1}{}}{}{\ {\normalfont\ttfamily ##1} :}\ %
    \{\ensuremath{##2}\}%
    {\normalfont\bfseries\,\par}%
  }
  \newcommand{\post}[2][]{%
    {\normalfont\bfseries\ttfamily Post}%
    \ifthenelse{\equal{##1}{}}{}{\ {\normalfont\ttfamily ##1} :}\
    \{\ensuremath{##2}\}%
    {\normalfont\bfseries\,\par}%
  }
  \renewcommand{\aux}[4]{%
    {\normalfont\bfseries\ttfamily aux\ }%
    {\normalfont\ttfamily ##1}%
    \ifthenelse{\equal{##2}{}}{}{\ (##2)}\ : ##3\, = \ensuremath{##4}%
    {\normalfont\bfseries\,;\par}%
  }
  \renewcommand{\pred}[3]{%
    {\normalfont\bfseries\ttfamily pred }%
    {\normalfont\ttfamily ##1}%
    \ifthenelse{\equal{##2}{}}{}{\ (##2) }%
    \{%
    \begin{adjustwidth}{+5em}{}
      \ensuremath{##3}
    \end{adjustwidth}
    \}%
    {\normalfont\bfseries\,\par}%
  }

  \newcommand{\res}{#1}
  \vspace{1ex}
  \noindent
  \encabezadoDeProc{#1}{#2}{#3}{#4}
  % Abrimos la llave
  \{\par%
  \tocarEspacios
}
% Ahora viene el cierre del ambiente problema
{
  % Cerramos la llave
  \noindent\}
  \vspace{1ex}
}


\newcommand{\aux}[4]{%
    {\normalfont\bfseries\ttfamily\noindent aux\ }%
    {\normalfont\ttfamily #1}%
    \ifthenelse{\equal{#2}{}}{}{\ (#2)}\ : #3\, = \ensuremath{#4}%
    {\normalfont\bfseries\,;\par}%
}


% \newcommand{\pre}[1]{\textsf{pre}\ensuremath{(#1)}}

\newcommand{\procnom}[1]{\textsf{#1}}
\newcommand{\procil}[3]{\textsf{proc #1}\ensuremath{(#2) = #3}}
\newcommand{\procilsinres}[2]{\textsf{proc #1}\ensuremath{(#2)}}
\newcommand{\preil}[2]{\textsf{Pre #1: }\ensuremath{#2}}
\newcommand{\postil}[2]{\textsf{Post #1: }\ensuremath{#2}}
\newcommand{\auxil}[2]{\textsf{fun }\ensuremath{#1 = #2}}
\newcommand{\auxilc}[4]{\textsf{fun }\ensuremath{#1( #2 ): #3 = #4}}
\newcommand{\auxnom}[1]{\textsf{fun }\ensuremath{#1}}
\newcommand{\auxpred}[3]{\textsf{pred }\ensuremath{#1( #2 ) \{ #3 \}}}

\newcommand{\comentario}[1]{{/*\ #1\ */}}

\newcommand{\nom}[1]{\ensuremath{\mathsf{#1}}}


% En las practicas/parciales usamos numeros arabigos para los ejercicios.
% Aca cambiamos los enumerate comunes para que usen letras y numeros
% romanos
\newcommand{\arreglarincisos}{%
  \renewcommand{\theenumi}{\alph{enumi}}
  \renewcommand{\theenumii}{\roman{enumii}}
  \renewcommand{\labelenumi}{\theenumi)}
  \renewcommand{\labelenumii}{\theenumii)}
}



%%%%%%%%%%%%%%%%%%%%%%%%%%%%%% PARCIAL %%%%%%%%%%%%%%%%%%%%%%%%
\let\@xa\expandafter
\newcommand{\tituloparcial}{\centerline{\depto -- \lamateria}
  \centerline{\elnombre -- \lafecha}%
  \setlength{\TPHorizModule}{10mm} % Fija las unidades de textpos
  \setlength{\TPVertModule}{\TPHorizModule} % Fija las unidades de
                                % textpos
  \arreglarincisos
  \newcounter{total}% Este contador va a guardar cuantos incisos hay
                    % en el parcial. Si un ejercicio no tiene incisos,
                    % cuenta como un inciso.
  \newcounter{contgrilla} % Para hacer ciclos
  \newcounter{columnainicial} % Se van a usar para los cline cuando un
  \newcounter{columnafinal}   % ejercicio tenga incisos.
  \newcommand{\primerafila}{}
  \newcommand{\segundafila}{}
  \newcommand{\rayitas}{} % Esto va a guardar los \cline de los
                          % ejercicios con incisos, asi queda mas bonito
  \newcommand{\anchodegrilla}{20} % Es para textpos
  \newcommand{\izquierda}{7} % Estos dos le dicen a textpos donde colocar
  \newcommand{\abajo}{2}     % la grilla
  \newcommand{\anchodecasilla}{0.4cm}
  \setcounter{columnainicial}{1}
  \setcounter{total}{0}
  \newcounter{ejercicio}
  \setcounter{ejercicio}{0}
  \renewenvironment{ejercicio}[1]
  {%
    \stepcounter{ejercicio}\textbf{\noindent Ejercicio \theejercicio. [##1
      puntos]}% Formato
    \renewcommand\@currentlabel{\theejercicio}% Esto es para las
                                % referencias
    \newcommand{\invariante}[2]{%
      {\normalfont\bfseries\ttfamily invariante}%
      \ ####1\hspace{1em}####2%
    }%
    \newcommand{\Proc}[5][result]{
      \encabezadoDeProc{####1}{####2}{####3}{####4}\hspace{1em}####5}%
  }% Aca se termina el principio del ejercicio
  {% Ahora viene el final
    % Esto suma la cantidad de incisos o 1 si no hubo ninguno
    \ifthenelse{\equal{\value{enumi}}{0}}
    {\addtocounter{total}{1}}
    {\addtocounter{total}{\value{enumi}}}
    \ifthenelse{\equal{\value{ejercicio}}{1}}{}
    {
      \g@addto@macro\primerafila{&} % Si no estoy en el primer ej.
      \g@addto@macro\segundafila{&}
    }
    \ifthenelse{\equal{\value{enumi}}{0}}
    {% No tiene incisos
      \g@addto@macro\primerafila{\multicolumn{1}{|c|}}
      \bgroup% avoid overwriting somebody else's value of \tmp@a
      \protected@edef\tmp@a{\theejercicio}% expand as far as we can
      \@xa\g@addto@macro\@xa\primerafila\@xa{\tmp@a}%
      \egroup% restore old value of \tmp@a, effect of \g@addto.. is

      \stepcounter{columnainicial}
    }
    {% Tiene incisos
      % Primero ponemos el encabezado
      \g@addto@macro\primerafila{\multicolumn}% Ahora el numero de items
      \bgroup% avoid overwriting somebody else's value of \tmp@a
      \protected@edef\tmp@a{\arabic{enumi}}% expand as far as we can
      \@xa\g@addto@macro\@xa\primerafila\@xa{\tmp@a}%
      \egroup% restore old value of \tmp@a, effect of \g@addto.. is
      % global
      % Ahora el formato
      \g@addto@macro\primerafila{{|c|}}%
      % Ahora el numero de ejercicio
      \bgroup% avoid overwriting somebody else's value of \tmp@a
      \protected@edef\tmp@a{\theejercicio}% expand as far as we can
      \@xa\g@addto@macro\@xa\primerafila\@xa{\tmp@a}%
      \egroup% restore old value of \tmp@a, effect of \g@addto.. is
      % global
      % Ahora armamos la segunda fila
      \g@addto@macro\segundafila{\multicolumn{1}{|c|}{a}}%
      \setcounter{contgrilla}{1}
      \whiledo{\value{contgrilla}<\value{enumi}}
      {%
        \stepcounter{contgrilla}
        \g@addto@macro\segundafila{&\multicolumn{1}{|c|}}
        \bgroup% avoid overwriting somebody else's value of \tmp@a
        \protected@edef\tmp@a{\alph{contgrilla}}% expand as far as we can
        \@xa\g@addto@macro\@xa\segundafila\@xa{\tmp@a}%
        \egroup% restore old value of \tmp@a, effect of \g@addto.. is
        % global
      }
      % Ahora armo las rayitas
      \setcounter{columnafinal}{\value{columnainicial}}
      \addtocounter{columnafinal}{-1}
      \addtocounter{columnafinal}{\value{enumi}}
      \bgroup% avoid overwriting somebody else's value of \tmp@a
      \protected@edef\tmp@a{\noexpand\cline{%
          \thecolumnainicial-\thecolumnafinal}}%
      \@xa\g@addto@macro\@xa\rayitas\@xa{\tmp@a}%
      \egroup% restore old value of \tmp@a, effect of \g@addto.. is
      \setcounter{columnainicial}{\value{columnafinal}}
      \stepcounter{columnainicial}
    }
    \setcounter{enumi}{0}%
    \vspace{0.2cm}%
  }%
  \newcommand{\tercerafila}{}
  \newcommand{\armartercerafila}{
    \setcounter{contgrilla}{1}
    \whiledo{\value{contgrilla}<\value{total}}
    {\stepcounter{contgrilla}\g@addto@macro\tercerafila{&}}
  }
  \newcommand{\grilla}{%
    \g@addto@macro\primerafila{&\textbf{TOTAL}}
    \g@addto@macro\segundafila{&}
    \g@addto@macro\tercerafila{&}
    \armartercerafila
    \ifthenelse{\equal{\value{total}}{\value{ejercicio}}}
    {% No hubo incisos
      \begin{textblock}{\anchodegrilla}(\izquierda,\abajo)
        \begin{tabular}{|*{\value{total}}{p{\anchodecasilla}|}c|}
          \hline
          \primerafila\\
          \hline
          \tercerafila\\
          \tercerafila\\
          \hline
        \end{tabular}
      \end{textblock}
    }
    {% Hubo incisos
      \begin{textblock}{\anchodegrilla}(\izquierda,\abajo)
        \begin{tabular}{|*{\value{total}}{p{\anchodecasilla}|}c|}
          \hline
          \primerafila\\
          \rayitas
          \segundafila\\
          \hline
          \tercerafila\\
          \tercerafila\\
          \hline
        \end{tabular}
      \end{textblock}
    }
  }%
  \vspace{0.4cm}
  \textbf{Nro. de orden:}

  \textbf{LU:}

  \textbf{Apellidos:}

  \textbf{Nombres:}

  \textbf{Nro. de hojas que adjunta:}
  \vspace{0.5cm}
}



% AMBIENTE CONSIGNAS
% Se usa en el TP para ir agregando las cosas que tienen que resolver
% los alumnos.
% Dentro del ambiente hay que usar \item para cada consigna

\newcounter{consigna}
\setcounter{consigna}{0}

\newenvironment{consignas}{%
  \newcommand{\consigna}{\stepcounter{consigna}\textbf{\theconsigna.}}%
  \renewcommand{\ejercicio}[1]{\item ##1 }
  \renewcommand{\proc}[5][result]{\item
    \encabezadoDeProc{##1}{##2}{##3}{##4}\hspace{1em}##5}%
  \newcommand{\invariante}[2]{\item%
    {\normalfont\bfseries\ttfamily invariante}%
    \ ##1\hspace{1em}##2%
  }
  \renewcommand{\aux}[4]{\item%
    {\normalfont\bfseries\ttfamily aux\ }%
    {\normalfont\ttfamily ##1}%
    \ifthenelse{\equal{##2}{}}{}{\ (##2)}\ : ##3 \hspace{1em}##4%
  }
  % Comienza la lista de consignas
  \begin{list}{\consigna}{%
      \setlength{\itemsep}{0.5em}%
      \setlength{\parsep}{0cm}%
    }
}%
{\end{list}}



% para decidir si usar && o ^
\newcommand{\y}[0]{\ensuremath{\land}}

% macros de correctitud
\newcommand{\semanticComment}[2]{#1 \ensuremath{#2};}
\newcommand{\namedSemanticComment}[3]{#1 #2: \ensuremath{#3};}


\newcommand{\local}[1]{\semanticComment{local}{#1}}

\newcommand{\vale}[1]{\semanticComment{vale}{#1}}
\newcommand{\valeN}[2]{\namedSemanticComment{vale}{#1}{#2}}
\newcommand{\impl}[1]{\semanticComment{implica}{#1}}
\newcommand{\implN}[2]{\namedSemanticComment{implica}{#1}{#2}}
\newcommand{\estado}[1]{\semanticComment{estado}{#1}}

\newcommand{\invarianteCN}[2]{\namedSemanticComment{invariante}{#1}{#2}}
\newcommand{\invarianteC}[1]{\semanticComment{invariante}{#1}}
\newcommand{\varianteCN}[2]{\namedSemanticComment{variante}{#1}{#2}}
\newcommand{\varianteC}[1]{\semanticComment{variante}{#1}}


\begin{document}

%\practica{4}{Precondición más débil en SmallLang}
\includepdf[pages=-]{practica5.pdf}

\begin{center}
\section*{RESOLUCIONES.}
\end{center}

\paragraph{Ejercicio 1}
\begin{enumerate}[label=\alph*)]
\item
	\begin{align*}
		P_c&: \{result=0\wedge i=0\}\\
		Q_c&: \{result=\sum_{j=0}^{|s|-1}s[j]\}
	\end{align*}
\item Si le sacas el igual a la segunda desigualdad no te agarra el ultimo elemento la 
	sumatoria.
\item Si le sacas el -1 al supraindice de la sumatoria entonces va a tratar de sumar la posicion del tamaño
	de la secuencia, como la secuencia no posee esta posicion se va a indefinir el resultado.
\item Si invertis de lugar las instrucciones del cuerpo del ciclo entonces no te agarra el primer elemento de la 
secuencia.
\item 
	Vamos  a demostrar los primeros 3 puntos del teorema del invariante. Esto nos permite afirmar que el ciclo, si termina entonces
	es correcto respecto a su especificacón. Especificacion del ciclo:
	\begin{itemize}
		\item $P_c: result=0 \wedge i=0$
		\item $Q_c:r esult=\sum_{j=0}^{|s|-1}s[j] $
		\item $I\equiv 0\leq i\leq |s| \wedge_L result=\sum_{j=0}^
			{i-1}s[j]$
	\end{itemize}
	\begin{flushleft}$\scalebox{1.3}{$P_c\Rightarrow I$}$\end{flushleft}
	Tenemos que vale $P_c$ como hipótesis. Queremos probar que vale I. Vamos a probarlo por partes:
	\begin{itemize}
		\item Queremos ver que vale $ 0\leq i\leq |s|$\smallskip \\
		Sabemos que i=0 (es información que nos da $P_c$). Entonces reemplazando, lo que queremos ver es que 
		$0\leq 0\leq |s|$. La primera parte, $0\leq 0$, es una tautología. Nos resta ver si $0\leq |s|$. Pero si la lista es
		vacia entonces $0\leq 0$, caso contrario si la lista fuera no vacia entonces $|s|>0$ y la desigualdad $0\leq |s|$ 
		se cumple.
		\item Queremos ver que $result=\sum_{j=0}^{i-1}s[j]$\\
		Como $i=0$, entonces $\sum_{j=0}^{i-1}s[j]=0=result$.
	\end{itemize}	 
	$\scalebox{1.3}{$(I\wedge \neg B)\Rightarrow Q_c$}$\smallskip \\
	Queremos demostrar que vale $Q_c$, asumiendo que valen tanto $I$ como $\neg B$.\\ Es decir, queremos probar que 
	$result =\sum_{j=0}^{|s|-1}s[j]$\\
	Como sabemos que vale $I$, podemos afirmar que $0\leq i\leq |s|$.\\ Además sabemos que vale $\neg B\equiv i\geq |s|$.
	\\ Luego $|s|\leq i\leq |s|$, entonces el unico valor que cumple esta condición es $i=|s|$. Analicemoslo.\\
	\begin{itemize}
	\item $i=|s|$
	\end{itemize}
	En este caso, podemos reemplazar $i$ por $|s|$ en la sumatoria del invariante y llegamos a $result =\sum_{j=0}^{|s|-1}s[j]$, exactamante lo que queríamos probar.\medskip \\
	$\scalebox{1.3}{$\{I\wedge B\} \textbf{ ciclo } \{I\}$}$\medskip \\
	Queremos ver que vale la siguiente tripla de Hoare $\{I\wedge B\} \textbf{ ciclo } \{I\}$.\\
	Llamemos S1 a la primer instrucción del cuerpo del ciclo, S2 a la segunda:\\
	S1$: result:=result+s[i];$\\
	S2$: i:=i+1$\\
	Lo primero que haremos es calcular $wp(ciclo,I)$.
	\begin{equation}wp(S1;S2,I)\stackrel{Ax3}{\equiv}wp(S1,wp(S2,I))\end{equation}
	Antes de seguir, debemos calcular $wp(S2,I)$. Para eso usaremos el axioma 1:
	\begin{align*}
	wp(S2,I)&\stackrel{Ax1}{\equiv}def(i+1)\wedge_L I_{i+1}^{i}\\
		&true\wedge_L (0\leq i+1\leq |s| \wedge_L result=\sum_{j=0}^{i+1-1}s[j])\\
		& (0\leq i+1\leq |s| \wedge_L result=\sum_{j=0}^{i+1-1}s[j])
	\end{align*}
	Volviendo a (1), reemplazamos $wp(S2,I)$ y nos queda:
	\begin{align*}
	wp(S1;S2,I)&\stackrel{Ax3}{\equiv}wp(S1,wp(S2,I))\equiv wp(S1,0\leq i+1\leq |s| \wedge_L result=\sum_{j=0}^{i+1-1}s[j])\\
		&\stackrel{Ax1}{\equiv}def(result+s[i])\wedge_L 0\leq i+1\leq |s| \wedge_L result+s[i]=\sum_{j=0}^{i}s[j]\\
		&\equiv 0\leq i<|s|\wedge_L 0\leq i+1\leq |s| \wedge_L result+s[i]=\sum_{j=0}^{i}s[j]\\
		&\equiv 0\leq i<|s| \wedge_L result=\sum_{j=0}^{i}s[j]-s[i]\\
		&\equiv 0\leq i<|s| \wedge_L result=\sum_{j=0}^{i-1}s[j]
	\end{align*}
	Una vez calculada la precondición más débil, debemos ver si $(I\wedge B)$ implican dicha precondición. Probamos cada
	parte por separado:
	\begin{itemize}
		\item $ 0\leq i< |s|$\smallskip \\
		Sabemos por $I$ que $0\leq i$, entonces la primera parte ya esta demostrada, ahora la interseccion entre la guarda
		y el invariante da como cota superior $i<|s|$, que es lo mismo que tengo en mi $wp$.
		\item $result=\sum_{j=0}^{i-1}s[j]$\smallskip \\
		Este $result$ es igual al del $I$.
		Como $(I\wedge B) \Rightarrow wp(ciclo,I)$, podemos afirmar que el cuerpo del ciclo preserva el invariante.
	\end{itemize}
	\item 
	Vamos  a demostrar los ultimos 2 puntos del teorema del invariante. Esto nos permite afirmar que el ciclo termina. Especificación del ciclo:
	\begin{itemize}
	\item $f_v:|s|-i$
	\end{itemize}
	$\scalebox{1.3}{$\{(I\wedge B\wedge v_0=f_v)\}\textbf{ ciclo }\{(f_v<v_0)\}$}$\medskip \\
	Dado que queremos demostrar que vale una tripla de Hoare, comenzaremos calculando la\\ precondición más débil $wp(ciclo,f_v<v_0)$.
	\begin{align*}
		wp(S1;S2,f_v<v_0)&\stackrel{Ax3}{\equiv}wp(S1,wp(S2,|s|-i<v_0))\\
			&\stackrel{Ax1}{\equiv}wp(S1,true\wedge_L |s|-(i+1)<v_0)\\
			&\stackrel{Ax3}{\equiv}true\wedge_L(true\wedge_L  |s|-(i+1)<v_0)\\
			&\equiv |s|-i-1<v_0
	\end{align*}
	Es decir, $wp(S1;S2,f_v<v_0)=|s|-i-1<v_0$. Ahora debemos ver que $(I\wedge B\wedge v_0=f_v)$ implican dicha WP. Parte de la 
	hipótesis es que $v_0=f_v$, es decir $v_0=|s|-i$. Restando 1 a ambos lados, $|s|-1-i=v_0-1<v_0$.\medskip \\
	$\scalebox{1.3}{$(I\wedge f_v\leq 0)\Rightarrow \neg B$}$\medskip\\
	Debemos mostrar que vale $\neg B$, es decir $i\geq |s|$.\\
	Sabemos que $f_v\leq 0$, es decir $|s|-i\leq 0$, luego $|s|\leq i$, como queriamos demostrar.
\end{enumerate}





\paragraph{Ejercicio 2}

Vamos  a demostrar los 5 puntos del teorema del invariante. Esto nos permite afirmar que el ciclo termina y es correcto
	 respecto a su especificacón (no demuestra que el programa entero sea correcto!). Especificacion del ciclo:
	\begin{itemize}
		\item $P_c: n\geq 0\wedge result=0 \wedge i=0$
		\item $Q_c:result=\sum_{j=0}^{n-1} $(if $j$ mod 2=0 then $j$ else 0 fi)
		\item $I\equiv 0\leq i\leq n+1\wedge i \textrm{ mod }2=0 \wedge result=\sum_{j=0}^{i-1}$(if $j$ mod 2=0 then $j$ else 0 fi)
		\item $f_v:n-i$
	\end{itemize}
	\begin{flushleft}$\scalebox{1.3}{$P_c\Rightarrow I$}$\end{flushleft}
	Tenemos que vale $P_c$ como hipótesis. Queremos probar que vale I. Vamos a probarlo por partes:
	\begin{itemize}
		\item Queremos ver que vale $ 0\leq i\leq n+1$\smallskip \\
		Sabemos que i=0 (es información que nos da $P_c$). Entonces reemplazando, lo que queremos ver es que 
		$0\leq 0\leq n+1$. La primera parte, $0\leq 0$, es una tautología. Nos resta ver si $0\leq n+1$. Pero $P_c$
		indica también que $n\leq 0$, luego $n+i\leq 1\leq 0$.
		\item Queremos ver que vale$i$ mod $2=0$\\
		Sabemos que $i=0$,luego podemos afirmar que $i$ mod $2=0$.
		\item Queremos ver que $result=\sum_{j=0}^{i-1}$(if $j$ mod 2=0 then $j$ else 0 fi)\\
		Como $i=0$, $result=\sum_{j=0}^{i-1}$(if $j$ mod 2=0 then $j$ else 0 fi)=$result=\sum_{j=0}^{-1}$(if $j$ mod 2=0 then $j$ else 0 fi).
		 Recordemos que si el rango de una sumatoria es vacío (como en este caso), la sumatoria tiene valor 0. Luego 
		 $result=\sum_{j=0}^{i-1}$(if $j$ mod 2=0 then $j$ else 0 fi)=0. Pero además sabemos que result = 0 (por $P_c$), asi que podemos
		  afirmar que 0=result=$result=\sum_{j=0}^{i-1}$(if $j$ mod 2=0 then $j$ else 0 fi)=0
	\end{itemize}	 
	$\scalebox{1.3}{$(I\wedge \neg B)\Rightarrow Q_c$}$\smallskip \\
	Queremos demostrar que vale $Q_c$, asumiendo que valen tanto $I$ como $\neg B$.\\ Es decir, queremos probar que 
	$result=\sum_{j=0}^{n-1}$(if $j$ mod 2=0 then $j$ else 0 fi)\\
	Como sabemos que vale $I$, podemos afirmar que $0\leq i\leq n+1$.\\ Además sabemos que vale $\neg B\equiv i\geq n$.
	\\ Luego $n\leq i\leq n+1$. Hay dos valores de $i$ que cumplen esa condición. Analicemos ambos casos:\\
	\begin{itemize}
	\item $i=n$
	\end{itemize}
	En este caso, podemos reemplazar $i$ por $n$ en la parte de la sumatoria del invariante y llegamos a
	 $result=\sum_{j=0}^{n-1}$(if $j$ mod 2=0 then $j$ else 0 fi), exactamante lo que queríamos probar.
	 \begin{itemize}
	\item $i=n+1$
	\end{itemize}
	En este caso, su hacemos el reemplazo, llegamos a  
	$result=\sum_{j=0}^{n}$(if $j$ mod 2=0 then $j$ else 0 fi) (y esto no es a lo que queremos llegar!).\\
	$result=\sum_{j=0}^{n}$(if $j$ mod 2=0 then $j$ else 0 fi) =$result=\sum_{j=0}^{n-1}$(if $j$ mod 2=0 then $j$ else 0 fi) 
	+(if $n$ mod 2=0 then $j$ else 0 fi). Sabemos además que $i$ mod 2=0 (información del invariante). Pero estamos en el caso en el cual
	 $i=n+1$. Entonces podemos afirmar que si $i$ es par, $n$ es impar. Luego $n$ no cumple la guada del IF y podemos afirmar que 
	 (if $n$ mod 2=0 then $j$ else 0 fi) = 0.\\
	 $result=\sum_{j=0}^{n}$(if $j$ mod 2=0 then $j$ else 0 fi) =$result=\sum_{j=0}^{n-1}$(if $j$ mod 2=0 then $j$ else 0 fi) +0.\medskip\\
	$\scalebox{1.3}{$\{I\wedge B\} \textbf{ ciclo } \{I\}$}$\medskip \\
	Queremos ver que vale la siguiente tripla de Hoare $\{I\wedge B\} \textbf{ ciclo } \{I\}$.\\
	Llamemos S1 a la primer instrucción del cuerpo del ciclo, S2 a la segunda:\\
	S1$: result:=result+i;$\\
	S2$: i:=i+2$\\
	Lo primero que haremos es calcular $wp(ciclo,I)$.
	\begin{equation}wp(S1;S2,I)\stackrel{Ax3}{\equiv}wp(S1,wp(S2,I))\end{equation}
	Antes de seguir, debemos calcular $wp(S2,I)$. Para eso usaremos el axioma 1:
	\begin{align*}
	wp(S2,I)&\stackrel{Ax1}{\equiv}def(i+2)\wedge_L I_{i+2}^{i}\\
		&\equiv true\wedge_L (0\leq i+2\leq n+1\wedge i+2\textrm{ mod }2=0 \wedge_L result=\sum_{j=0}^{i+2-1}\textrm{(if j mod 2=0 then j else 0 fi)})\\
		&\equiv (0\leq i+2\leq n+1\wedge i+2\textrm{ mod }2=0 \wedge_L result=\sum_{j=0}^{i+1}\textrm{(if j mod 2=0 then j else 0 fi)})
	\end{align*}
	Volviendo a (2), reemplazamos $wp(S2,I)$ y nos queda:
	\begin{align*}
	wp(S1;S2,I)&\stackrel{Ax3}{\equiv}wp(S1,wp(S2,I))\equiv wp(S1,0\leq i+2\leq n+1 \wedge i+2\textrm{ mod }2=0 
			\wedge result=\sum_{j=0}^{i+1}\textrm{(if j mod 2=0 then j else 0 fi)})\\
		&\stackrel{Ax1}{\equiv}def(result+i)\wedge_L (0\leq i+2\leq n+1 \wedge i+2\textrm{ mod }2=0
			 \wedge result=\sum_{j=0}^{i+1}\textrm{(if j mod 2=0 then j else 0 fi)})\\
		&\equiv true \wedge_L (0\leq i+2\leq n+1 \wedge i+2\textrm{ mod }2=0
			 \wedge result=\sum_{j=0}^{i+1}\textrm{(if j mod 2=0 then j else 0 fi)})\\
	\end{align*}
	Una vez calculada la precondición más débil, debemos ver si $(I\wedge B)$ implican dicha precondición. Probamos cada
	parte por separado:
	\begin{itemize}
		\item $ 0\leq i+2\leq n+1$\smallskip \\
		Sabemos por $I$ que $ i>0$, luego podemos afirmar que $0\leq i+2$\\
		Sabemos por $B$ que $i<n$, luego (sumando 2 en ambos términos): $i+2<n+2$, lo cual es equivalente a decir
		que $i+2\leq n+1$
		\item $i+2$ mod 2=0\smallskip \\
		Sabemos por $I$ que $i$ mod 2=0. Si $i$ es par, al sumarle 2 sigue siendo par, luego i +2 mod 2=0 vale. 
		\item $result+i=\sum_{j=0}^{i-1}\textrm{(if j mod 2=0 then j else 0 fi)}$\smallskip \\
		La sumatoria puede separarse en 3 términos:\\
		$\sum_{j=0}^{i-1}\textrm{(if j mod 2=0 then j else 0 fi)}+$\\
		if i mod 2=0 then i else 0 fi\\
		if i+1 mod 2=0 then i+1 else 0 fi\\
		El primero de los 3 términos es igual a $result$ (lo sabemos por $I$).\\
		El segundo término es $i$ (ya que por $I$ sabemos que $i$ es par).\\
		El tercer término es 0 (ya que por $I$ sabemos que $i$ es par, y por lo tanto $i+1$ es impar).\\
		Entonces sumando los 3 términos nos queda: $result +i+0$, que es lo que esperábamos que valiera la sumatoria 
		$\sum_{j=0}^{i+1}(...)$\\
		Como $(I\wedge B)\Rightarrow wp(ciclo,I)$,podemos afirmar que el cuerpo del ciclo preserva el invariante.
	\end{itemize}
	$\scalebox{1.3}{$\{(I\wedge B\wedge v_0=f_v)\}\textbf{ ciclo }\{(f_v<v_0)\}$}$\medskip \\
	Dado que queremos demostrar que vale una tripla de Hoare, comenzaremos calculando la\\ precondición más débil $wp(ciclo,f_v<v_0)$.
	\begin{align*}
		wp(S1;S2,f_v<v_0)&\stackrel{Ax3}{\equiv}wp(S1,wp(S2,n-i<v_0))\\
			&\stackrel{Ax1}{\equiv}wp(S1,true\wedge_L n-(i+2)<v_0)\\
			&\stackrel{Ax3}{\equiv}true\wedge_L(true\wedge_L  n-(i+2)<v_0)\\
			&\equiv n-i-2<v_0
	\end{align*}
	Es decir, $wp(S1;S2,f_v<v_0)=n-i-2<v_0$. Ahora debemos ver que $(I\wedge B\wedge v_0=f_v)$ implican dicha WP. Parte de la 
	hipótesis es que $v_0=f_v$, es decir $v_0=n-i$. Restando 2 a ambos lados, $n-i-2=v_0-2<v_0$.\medskip \\
	$\scalebox{1.3}{$(I\wedge f_v\leq 0)\Rightarrow \neg B$}$\medskip\\
	Debemos mostrar que vale $\neg B$, es decir $i\geq n$.\\
	Sabemos que $f_v\leq 0$, es decir $n-i\leq 0$, luego $n\leq i$, como queriamos demostrar.
	
	
	
	
	
\paragraph{Ejercicio 3\\}
\noindent
\textbf{proc exponenciacion}(in $s: \mathbb{Z}$,in $n:\mathbb{Z}$, out $result: \mathbb{Z})\ \{$\smallskip \\
\hspace*{6mm}\textbf{Pre}$\{ n\geq 0 \wedge \neg(m=0\wedge n=0)\}$\\
\hspace*{6mm}\textbf{Post}$\{ result=m^n\}$\\
$\}$\medskip\\

\noindent
\begin{enumerate}[label=\alph*)]
	\item 

\begin{lstlisting}
i:=1;
result:=m;

|\textbf{while}|(i<n) |\textbf{do}|
  result=:result*m;
  i:=i+1
|\textbf{endwhile}|
\end{lstlisting}

 Especificacion del ciclo:
	\begin{itemize}
		\item $P_c: result=m \wedge i=1\wedge n\geq 0 \wedge \neg(m=0\wedge n=0)$
		\item $Q_c:result=m^n$
		\item $I\equiv 1\leq i\leq n \wedge result=m^i$
		\item $f_v:n-i$
	\end{itemize}
	Solo demostraremos correción parcial de este ciclo.
	\begin{flushleft}$\scalebox{1.3}{$P_c\Rightarrow I$}$\end{flushleft}
	Debemos demostrar que vale $I$ sabiendo que vale $P_c$ 
	\begin{itemize}
		\item $ 1\leq i\leq n$\smallskip \\
		Por $P_c$ sabemos que $i=1$, entonces, la primera desigualdad es una tautologia,
		$1\leq 1$, para la segunda vemos que $n\geq0\wedge \neg(m=0 \wedge n=0)$, luego como 
		$n$, no puede ser 0 y es mayor o igual a cero tengo que $1\leq n$.
		\item  $result=m^i$\\
		Trivial, por $P_c$ $result=m$, y como $i=1$, entonces $m^i=m$, por lo que 
		 la igualdad vale.  
	\end{itemize}	 
	$\scalebox{1.3}{$(I\wedge \neg B)\Rightarrow Q_c$}$\smallskip \\
	Queremos demostrar que vale $Q_c$, asumiendo que vale  $I\wedge \neg B$.
	\begin{itemize}
	\item cuando deja de valer la guarda, $i=n$ sabemos por el invariante que $result=m^i$, como $i=n$, entonces,
	vale que $result=m^n$, que es exactamente lo mismo que $Q_c$
	\end{itemize}
	$\scalebox{1.3}{$\{I\wedge B\} \textbf{ ciclo } \{I\}$}$\medskip \\
	Queremos ver que vale la siguiente tripla de Hoare $\{I\wedge B\} \textbf{ ciclo } \{I\}$.\\
	Llamemos S1 a la primer instrucción del cuerpo del ciclo, S2 a la segunda:\\
	S1$: result:=result*m;$\\
	S2$: i:=i+1$\\
	Lo primero que haremos es calcular $wp(ciclo,I)$.
	\begin{equation}wp(S1;S2,I)\stackrel{Ax3}{\equiv}wp(S1,wp(S2,I))\end{equation}
	Antes de seguir, debemos calcular $wp(S2,I)$. Para eso usaremos el axioma 1:
	\begin{align*}
	wp(S2,I)&\stackrel{Ax1}{\equiv}def(i+1)\wedge_L I_{i+1}^{i}\\
		&\equiv true\wedge_L (1\leq i+1\leq n \wedge_L result=m^{i+1})\\
		&\equiv (1\leq i+1\leq n \wedge_L result=m^{i+1})
	\end{align*}
	Volviendo a (3), reemplazamos $wp(S2,I)$ y nos queda:
	\begin{align*}
	wp(S1;S2,I)&\stackrel{Ax3}{\equiv}wp(S1,wp(S2,I))\equiv wp(S1,1\leq i+1\leq n \wedge_L result=m^{i+1})\\
		&\stackrel{Ax1}{\equiv}def(result*m)\wedge_L 1\leq i+1\leq n \wedge_L result*m=m^{i+1}\\
		&\equiv 1\leq i+1\leq n \wedge_L result*m=m^{i+1}
	\end{align*}
	Una vez calculada la precondición más débil, debemos ver si $(I\wedge B)$ implican dicha precondición. Probamos cada
	parte por separado:
	\begin{itemize}
		\item $ 0\leq i+1\leq n$\smallskip \\
		Sabemos por $I$ que $ i>0$, luego podemos afirmar que $1\leq i+1$\\
		Sabemos por $B$ que $i<n$, luego $i+1\geq n$.
		\item $result*m=m^{i+1}$\smallskip \\
		Divido en ambos miembros por $m$, y tengo lo mismo que en mi invariante.\smallskip \\
		Como $(I\wedge B) \Rightarrow wp(ciclo,I)$, podemos afirmar que el cuerpo del ciclo preserva el invariante.
	\end{itemize}
	$\scalebox{1.3}{$\{(I\wedge B\wedge v_0=f_v)\}\textbf{ ciclo }\{(f_v<v_0)\}$}$\medskip \\
	Dado que queremos demostrar que vale una tripla de Hoare, comenzaremos calculando la\\ precondición más débil $wp(ciclo,f_v<v_0)$.
	\begin{align*}
		wp(S1;S2,f_v<v_0)&\stackrel{Ax3}{\equiv}wp(S1,wp(S2,n-i<v_0))\\
			&\stackrel{Ax1}{\equiv}wp(S1,true\wedge_L n-(i+1)<v_0)\\
			&\stackrel{Ax3}{\equiv}true\wedge_L(true\wedge_L  n-(i+1)<v_0)\\
			&\equiv n-i-1<v_0
	\end{align*}
	Es decir, $wp(S1;S2,f_v<v_0)=n-i-1<v_0$. Ahora debemos ver que $(I\wedge B\wedge v_0=f_v)$ implican dicha WP. Parte de la 
	hipótesis es que $v_0=f_v$, es decir $v_0=n-i$. Restando 1 a ambos lados, $n-1-i=v_0-1<v_0$.\medskip \\
	$\scalebox{1.3}{$(I\wedge f_v\leq 0)\Rightarrow \neg B$}$\medskip\\
	Debemos mostrar que vale $\neg B$, es decir $i\geq n$.\\
	Sabemos que $f_v\leq 0$, es decir $n-i\leq 0$, luego $n\leq i$, como queriamos demostrar.
	
	\item falla el primero, el segundo el tercero el cuarto y el quinto, falla todo para ser más preciso.
	\item creo que no falla, empieza con el caso $m^0=1$, y sigue como en mi demostración en a), no lo
	voy a probar de vuelta.
	\item habria que agregar en la pre que $n\geq 2$.
\end{enumerate}

\paragraph{Ejercicio 4\\}
\noindent
\textbf{proc sumaDivisores}(in $n: \mathbb{Z}$, out $result: \mathbb{Z})\ \{$\smallskip \\
\hspace*{6mm}\textbf{Pre}$\{ n\geq 1 \}$\\
\hspace*{6mm}\textbf{Post}$\{ result=\sum_{j=1}^{n}(\textrm{if}\ n\ mod\ j=0\ \textrm{then}\ j\ \textrm{else}\ 0\ \textrm{fi})\}$\\
$\}$\medskip\\

\noindent
\begin{enumerate}[label=\alph*)]
	\item 

\begin{lstlisting}
i:=0;
result:=0;

|\textbf{while}|(i<n) |\textbf{do}|
  i:=i+1
  |\textbf{if}| (n mod i = 0) 
     result=:result +i
  |\textbf{else}|
     skip
  |\textbf{endif}|;
|\textbf{endwhile}|
\end{lstlisting}
	\item
	
 Especificacion del ciclo:
	\begin{itemize}
		\item $P_c: result=0 \wedge i=0\wedge n\geq 1$
		\item $Q_c:result=\sum_{j=1}^{n}(\textrm{if}\ n\ mod\ j=0\ \textrm{then}\ j\ \textrm{else}\ 0\ \textrm{fi})$
		\item $I\equiv 1\leq i\leq n \wedge result=\sum_{j=1}^{i}(\textrm{if}\ n\ mod\ j=0\ \textrm{then}\ j\ \textrm{else}\ 0\ \textrm{fi})$
		\item $f_v:n-i$
	\end{itemize}
	
	Solo demostraremos correción parcial de este ciclo.
	\begin{flushleft}$\scalebox{1.3}{$P_c\Rightarrow I$}$\end{flushleft}
	Debemos demostrar que vale $I$ sabiendo que vale $P_c$ 
	\begin{itemize}
		\item $ 1\leq i\leq n$\smallskip \\
		Por $P_c$ sabemos que $i=1$, entonces, la primera desigualdad es una tautologia,
		$1\leq 1$, para la segunda vemos que $n\geq 1$, luego como 
		$n$, no puede ser 0 y es mayor o igual a 1 tengo que $1\leq n$.
		\item  $result=\sum_{j=1}^{i}(\textrm{if}\ n\ mod\ j=0\ \textrm{then}\ j\ \textrm{else}\ 0\ \textrm{fi})$\\
		Trivial, por $P_c$ $result=0$, y  como $i=0$, entonces
		 $\sum_{j=1}^{0}(\textrm{if}\ n\ mod\ j=0\ \textrm{then}\ j\ \textrm{else}\ 0\ \textrm{fi})=result=0$.
	\end{itemize}	 
	$\scalebox{1.3}{$(I\wedge \neg B)\Rightarrow Q_c$}$\smallskip \\
	Queremos demostrar que vale $Q_c$, asumiendo que vale  $I\wedge \neg B$.
	\begin{itemize}
	\item cuando deja de valer la guarda, $i=n$ sabemos por el invariante que $result=\sum_{j=1}^{i}(\textrm{if}\ n\ mod\ j=0\ \textrm{then}\ j\ \textrm{else}\ 0\ \textrm{fi})$, como $i=n$, entonces,
	vale que $result=\sum_{j=1}^{n}(\textrm{if}\ n\ mod\ j=0\ \textrm{then}\ j\ \textrm{else}\ 0\ \textrm{fi})$, que es exactamente lo mismo que $Q_c$
	\end{itemize}
	$\scalebox{1.3}{$\{I\wedge B\} \textbf{ ciclo } \{I\}$}$\medskip \\
	Queremos ver que vale la siguiente tripla de Hoare $\{I\wedge B\} \textbf{ ciclo } \{I\}$.\\
	Llamemos S1 a la primer instrucción del cuerpo del ciclo, S2 a la segunda:\\
	S1$: i:=i+1$\\
	S2$: \textbf{if}\ (n\ mod\ i=0)\ \textbf{then}\ result:=result+i\ \textbf{else}\ skip\ \textbf{endif};$\\
	Lo primero que haremos es calcular $wp(ciclo,I)$.
	\begin{equation}wp(S1;S2,I)\stackrel{Ax3}{\equiv}wp(S1,wp(S2,I))\end{equation}
	Antes de seguir, debemos calcular $wp(S2,I)$. Para eso usaremos el axioma 1:
	\begin{align*}
	wp(S2,I)&\stackrel{Ax4}{\equiv}def(n\ mod\ i=0)\wedge_L \Bigg(\Big((n\ mod\ i=0)\wedge (wp(result:=result+i,I))\Big) 
   				\vee \Big(\neg  (n\ mod\ i=0)\wedge (wp(skip,I))\Big)\Bigg)\\
		&\equiv True \wedge_L \Bigg(\Big((n\ mod\ i=0)\wedge (def(result+i)\wedge_LI_{result+i}^{result}))\Big) 
   				\vee \Big(\neg  (n\ mod\ i=0)\wedge (I)\Big)\Bigg)\\
		&\equiv \Bigg(\Big((n\ mod\ i=0)\wedge (I_{result+i}^{result}))\Big) 
   				\vee \Big(\neg  (n\ mod\ i=0)\wedge (I)\Big)\Bigg)\\
	\end{align*}
	Volviendo a (4), reemplazamos $wp(S2,I)$ y nos queda:
	\begin{align*}
	wp(S1;S2,I)&\stackrel{Ax3}{\equiv}wp(S1,wp(S2,I))\\
	&\equiv  wp(S1,\Big((n\ mod\ i=0)\wedge (1\leq i\leq n \wedge result+1=\sum_{j=1}^{i}(\textrm{if}\ n\ mod\ j=0\ \textrm{then}\ j\ \textrm{else}\ 0\ \textrm{fi})))\Big) 
   				\\ &\vee \Big(\neg  (n\ mod\ i=0)\wedge (1\leq i\leq n \wedge result=\sum_{j=1}^{i}(\textrm{if}\ n\ mod\ j=0\ \textrm{then}\ j\ \textrm{else}\ 0\ \textrm{fi}))\Big))\\
		&\stackrel{Ax1}{\equiv}def(i+1)\wedge_L \Big((n\ mod\ i+1=0)\wedge (1\leq i+1\leq n \wedge result+1=\sum_{j=1}^{i+1}(\textrm{if}\ n\ mod\ j=0\ \textrm{then}\ j\ \textrm{else}\ 0\ \textrm{fi})))\Big) 
   				\\ &\vee \Big(\neg  (n\ mod\ i+1=0)\wedge (1\leq i+1\leq n \wedge result=\sum_{j=1}^{i+1}(\textrm{if}\ n\ mod\ j=0\ \textrm{then}\ j\ \textrm{else}\ 0\ \textrm{fi}))\Big))\\
		&\equiv  \Big((n\ mod\ i+1=0)\wedge (1\leq i+1\leq n \wedge result+1=\sum_{j=1}^{i+1}(\textrm{if}\ n\ mod\ j=0\ \textrm{then}\ j\ \textrm{else}\ 0\ \textrm{fi})))\Big) 
   				\\ &\vee \Big(\neg  (n\ mod\ i+1=0)\wedge (1\leq i+1\leq n \wedge result=\sum_{j=1}^{i+1}(\textrm{if}\ n\ mod\ j=0\ \textrm{then}\ j\ \textrm{else}\ 0\ \textrm{fi}))\Big))\\
	\end{align*}
	Una vez calculada la precondición más débil, debemos ver si $(I\wedge B)$ implican dicha precondición. Probamos cada
	parte por separado:
	\begin{itemize}
		\item $\Big((n\ mod\ i+1=0)\wedge (1\leq i+1\leq n \wedge result+1=\sum_{j=1}^{i+1}(\textrm{if}\ n\ mod\ j=0\ \textrm{then}\ j\ \textrm{else}\ 0\ \textrm{fi})))\Big) $\smallskip \\
		Ya lo hice en la guia 4
		\item $\Big(\neg  (n\ mod\ i+1=0)\wedge (1\leq i+1\leq n \wedge result=\sum_{j=1}^{i+1}(\textrm{if}\ n\ mod\ j=0\ \textrm{then}\ j\ \textrm{else}\ 0\ \textrm{fi}))\Big))$\smallskip \\
		Ya lo hice en la guia 4\smallskip \\
		Como $(I\wedge B) \Rightarrow wp(ciclo,I)$, podemos afirmar que el cuerpo del ciclo preserva el invariante.
	\end{itemize}
\end{enumerate}

\paragraph{Ejercicio 5\\}
\noindent
\textbf{proc sumaPosicionesImpares}(in $s: seq\langle\mathbb{Z}\rangle$, out $result: \mathbb{Z})\ \{$\smallskip \\
\hspace*{6mm}\textbf{Pre}$\{True \}$\\
\hspace*{6mm}\textbf{Post}$\{ result=\sum_{i=0}^{|s|-1}(\textrm{if}\ i\ mod\ 2=1\ \textrm{then}\ s[i]\ \textrm{else}\ 0\ \textrm{fi})\}$\\
$\}$\medskip\\

\noindent
\begin{enumerate}[label=\alph*)]
	\item 

\begin{lstlisting}
j:=0;
result:=0;

|\textbf{while}|(j<s.size()) |\textbf{do}|
  |\textbf{if}| (j mod 2 = 1) 
     result:=result +s[j]
  |\textbf{else}|
     skip
  |\textbf{endif}|;
   j:=j+1
|\textbf{endwhile}|
\end{lstlisting}
	\item
	
 Especificacion del ciclo:
	\begin{itemize}
		\item $P_c: result=0 \wedge j=0$
		\item $Q_c: result=\sum_{i=0}^{|s|-1}(\textrm{if}\ i\ mod\ 2=1\ \textrm{then}\ s[i]\ \textrm{else}\ 0\ \textrm{fi}))$
		\item $I\equiv 0\leq j\leq |s| \wedge result=\sum_{i=0}^{j-1}(\textrm{if}\ i\ mod\ 2=1\ \textrm{then}\ s[i]\ \textrm{else}\ 0\ \textrm{fi})$
		\item $f_v:n-i$
	\end{itemize}
	
	Solo demostraremos correción parcial de este ciclo.
	\begin{flushleft}$\scalebox{1.3}{$P_c\Rightarrow I$}$\end{flushleft}
	Debemos demostrar que vale $I$ sabiendo que vale $P_c$ 
	\begin{itemize}
		\item $ 0\leq j\leq |s|$\smallskip \\
		Por $P_c$ sabemos que $j=0$, entonces, la primera desigualdad es una tautologia,
		$0\leq 0$, para la segunda vemos que $|s|\geq 0$, luego como 
		$|s|$, no puede ser menor que cero es valido.
		\item  $result=\sum_{i=0}^{j-1}(\textrm{if}\ i\ mod\ 2=1\ \textrm{then}\ s[i]\ \textrm{else}\ 0\ \textrm{fi})$\\
		Trivial, por $P_c$ $result=0$, y  como $j=0$, entonces
		 $\sum_{i=0}^{-1}(\textrm{if}\ i\ mod\ 2=1\ \textrm{then}\ s[i]\ \textrm{else}\ 0\ \textrm{fi})=result=0$.
	\end{itemize}	 
	$\scalebox{1.3}{$(I\wedge \neg B)\Rightarrow Q_c$}$\smallskip \\
	Queremos demostrar que vale $Q_c$, asumiendo que vale  $I\wedge \neg B$.
	\begin{itemize}
	\item cuando deja de valer la guarda, $j=|s|$ sabemos por el invariante que $result=\sum_{i=0}^{j-1}(\textrm{if}\ i\ mod\ 2=1\ \textrm{then}\ s[i]\ \textrm{else}\ 0\ \textrm{fi})$, como $j=|s|$, entonces,
	vale que $result=\sum_{j=1}^{|s|-1}(\textrm{if}\ n\ mod\ j=0\ \textrm{then}\ j\ \textrm{else}\ 0\ \textrm{fi})$, que es exactamente lo mismo que $Q_c$
	\end{itemize}
	$\scalebox{1.3}{$\{I\wedge B\} \textbf{ ciclo } \{I\}$}$\medskip \\
	Queremos ver que vale la siguiente tripla de Hoare $\{I\wedge B\} \textbf{ ciclo } \{I\}$.\\
	Llamemos S1 a la primer instrucción del cuerpo del ciclo, S2 a la segunda:\\
	S1$: \textbf{if}\ (j\ mod\ 2=1\ \textbf{then}\ result:=result+s[j]\ \textbf{else}\ skip\ \textbf{endif};$\\
	S2$: j:=j+1$\\
	Lo primero que haremos es calcular $wp(ciclo,I)$.
	\begin{equation}wp(S1;S2,I)\stackrel{Ax3}{\equiv}wp(S1,wp(S2,I))\end{equation}
	Antes de seguir, debemos calcular $wp(S2,I)$. Para eso usaremos el axioma 1:
	\begin{align*}
	wp(S2,I)&\stackrel{Ax1}{\equiv}def(j+1)\wedge_L I_{j+1}^{j}\\
		&\equiv true\wedge_L 0\leq j+1\leq |s| \wedge result=\sum_{i=0}^{j-1+1}(\textrm{if}\ i\ mod\ 2=1\ \textrm{then}\ s[i]\ \textrm{else}\ 0\ \textrm{fi})\\
		&\equiv 0\leq j+1\leq |s| \wedge result=\sum_{i=0}^{j-1+1}(\textrm{if}\ i\ mod\ 2=1\ \textrm{then}\ s[i]\ \textrm{else}\ 0\ \textrm{fi}))
	\end{align*}
	
	Volviendo a (4), reemplazamos $wp(S2,I)$ y nos queda:
	\begin{align*}
	wp(S1;S2,I)&\stackrel{Ax3}{\equiv}wp(S1,wp(S2,I))\\
	&\equiv  wp(S1,0\leq j+1\leq |s| \wedge result=\sum_{i=0}^{j-1+1}(\textrm{if}\ i\ mod\ 2=1\ \textrm{then}\ s[i]\ \textrm{else}\ 0\ \textrm{fi}))\\
	&\equiv  wp(S1,I_{j+1}^{j})\\
	&\stackrel{Ax4}{\equiv}def(j\ mod\ 2=1)\wedge_L \Bigg(\Big((j\ mod\ 2=1)\wedge (wp(result:=result+s[j],I_{j+1}^{j}))\Big) 
   				\\&\vee \Big(\neg  (j\ mod\ 2=1)\wedge (wp(skip,I_{j+1}^{j}))\Big)\Bigg)\\
		&\stackrel{Ax1}{\equiv} True\wedge_L \Bigg(\Big((j\ mod\ 2=1)\wedge (wp(result:=result+s[j],I_{j+1}^{j}))\Big) 
   				\vee \Big(\neg  (j\ mod\ 2=1)\wedge I_{j+1}^{j})\Big)\Bigg)\\
		&\equiv  \Bigg(\Big((j\ mod\ 2=1)\wedge (wp(result:=result+s[j],I_{j+1}^{j}))\Big) 
   				\vee \Big(\neg  (j\ mod\ 2=1)\wedge I_{j+1}^{j})\Big)\Bigg)\\
   		&\equiv  \Bigg(\Big((j\ mod\ 2=1)\wedge (0\leq j+1\leq |s| \wedge result+s[j]=\sum_{i=0}^{j-1+1}
   			(\textrm{if}\ i\ mod\ 2=1\ \textrm{then}\ s[i]\ \textrm{else}\ 0\ \textrm{fi}))\Big) 
   				\\&\vee \Big(\neg  (j\ mod\ 2=1)\wedge 0\leq j+1\leq |s| \wedge result=\sum_{i=0}^{j-1+1}(\textrm{if}\ i\ mod\ 2=1\ \textrm{then}\ s[i]\ \textrm{else}\ 0\ \textrm{fi}))\Big)\Bigg)\\
	\end{align*}
	Una vez calculada la precondición más débil, debemos ver si $(I\wedge B)$ implican dicha precondición. Probamos cada
	parte por separado:
	\begin{itemize}
		\item $\Big((j\ mod\ 2=1)\wedge (0\leq j+1\leq |s| \wedge result+s[j]=\sum_{i=0}^{j-1+1}
   			(\textrm{if}\ i\ mod\ 2=1\ \textrm{then}\ s[i]\ \textrm{else}\ 0\ \textrm{fi}))\Big) $\smallskip \\
		Supongo que vale la izquierda y entonces vale la derecha.
		\item $\Big(\neg  (j\ mod\ 2=1)\wedge 0\leq j+1\leq |s| \wedge result=\sum_{i=0}^{j-1+1}(\textrm{if}\ i\ mod\ 2=1\ \textrm{then}\ s[i]\ \textrm{else}\ 0\ \textrm{fi}))\Big)$\smallskip \\
		Supongo que vale la izquierda y entonces vale la derecha.\smallskip \\
		Como $(I\wedge B) \Rightarrow wp(ciclo,I)$, podemos afirmar que el cuerpo del ciclo preserva el invariante.
	\end{itemize}
\end{enumerate}

\paragraph{Ejercicio 6\\}
\noindent
\textbf{proc maximo}(in $s: seq\langle\mathbb{Z}\rangle$, out $i: \mathbb{Z})\ \{$\smallskip \\
\hspace*{6mm}\textbf{Pre}$\{|s|\geq 1 \}$\\
\hspace*{6mm}\textbf{Post}$\{ 0\leq i<|s|\wedge_L (\forall j:\mathbb{Z})(0\leq j<|s|\rightarrow_L s[j]\leq s[i])\}$\\
$\}$\medskip\\

\noindent
\begin{enumerate}[label=\alph*)]
	\item 

\begin{lstlisting}
i:=0;
j:=1;

|\textbf{while}|(j<s.size()) |\textbf{do}|
  |\textbf{if}| (s[j] > s[i]) 
     i := j
  |\textbf{else}|
     skip
  |\textbf{endif}|;
   j:=j+1
|\textbf{endwhile}|
\end{lstlisting}
	\item
	
 Especificacion del ciclo:
	\begin{itemize}
		\item $P_c: i=0 \wedge j=1\wedge |s|\geq 1$
		\item $Q_c: 0\leq i<|s|\wedge_L (\forall j:\mathbb{Z})(0\leq j<|s|\rightarrow_L s[j]\leq s[i])$
		\item $I\equiv (0\leq i <|s|\wedge 1\leq j\leq |s|)\wedge_L (\forall k:\mathbb{Z})(0\leq k<j\rightarrow_L s[k]\leq s[i])$
		\item $f_v:|s|+1-j$
	\end{itemize}
	
	Solo demostraremos correción parcial de este ciclo.
	\begin{flushleft}$\scalebox{1.3}{$P_c\Rightarrow I$}$\end{flushleft}
	Debemos demostrar que vale $I$ sabiendo que vale $P_c$ 
	\begin{itemize}
		\item $ (0\leq i <|s|\wedge 1\leq j\leq |s|)$\smallskip \\
		Por $P_c$ sabemos que $j=1,i=0$, entonces, la primera desigualdad de ambos terminos es una tautologia,
		$0\leq 0\wedge 1\leq 1$, para la segunda vemos que por $P_c$ $|s|\geq 1$, luego como $i,j$ son o menores o iguales
		a 1 vale.
		\item  $(\forall k:\mathbb{Z})(0\leq k<j\rightarrow_L s[k]\leq s[i])$\\
		Parto desde que $i=0,j=1$
		El unico .$k$ que cumple que $k<1$ es 0, por lo que es cierto que $s[0]\leq s[0]$. 
	\end{itemize}	 
	$\scalebox{1.3}{$(I\wedge \neg B)\Rightarrow Q_c$}$\smallskip \\
	Queremos demostrar que vale $Q_c$, asumiendo que vale  $I\wedge \neg B$.
	\begin{itemize}
	\item cuando deja de valer la guarda, $j=|s|$ reemplazamos $j$ por $|s|$ en el invariante y tenemos $Q_c$
	\end{itemize}
	$\scalebox{1.3}{$\{I\wedge B\} \textbf{ ciclo } \{I\}$}$\medskip \\
	Queremos ver que vale la siguiente tripla de Hoare $\{I\wedge B\} \textbf{ ciclo } \{I\}$.\\
	Llamemos S1 a la primer instrucción del cuerpo del ciclo, S2 a la segunda:\\
	S1$: \textbf{if}\ (s[j]>s[i])\ \textbf{then}\ i:=j\ \textbf{else}\ skip\ \textbf{endif};$\\
	S2$: j:=j+1$\\
	Lo primero que haremos es calcular $wp(ciclo,I)$.
	\begin{equation}wp(S1;S2,I)\stackrel{Ax3}{\equiv}wp(S1,wp(S2,I))\end{equation}
	Antes de seguir, debemos calcular $wp(S2,I)$. Para eso usaremos el axioma 1:
	\begin{align*}
	wp(S2,I)&\stackrel{Ax1}{\equiv}def(j+1)\wedge_L I_{j+1}^{j}\\
		&\equiv true\wedge_L (0\leq i <|s|\wedge 1\leq j+1\leq |s|)\wedge_L (\forall k:\mathbb{Z})(0\leq k<j+1\rightarrow_L s[k]\leq s[i])\\
		&\equiv (0\leq i <|s|\wedge 1\leq j+1\leq |s|)\wedge_L (\forall k:\mathbb{Z})(0\leq k<j+1\rightarrow_L s[k]\leq s[i])
	\end{align*}
	
	Volviendo a (4), reemplazamos $wp(S2,I)$ y nos queda:
	\begin{align*}
	wp(S1;S2,I)&\stackrel{Ax3}{\equiv}wp(S1,wp(S2,I))\\
	&\equiv  wp(S1,I_{j+1}^{j})\\
	&\stackrel{Ax4}{\equiv}def(s[j+1]>s[i])\wedge_L \Bigg(\Big((s[j+1]>s[i])\wedge (wp(i:=j,I_{j+1}^{j}))\Big) 
   				\\&\vee \Big(\neg  (s[j+1]>s[i])\wedge (wp(skip,I_{j+1}^{j}))\Big)\Bigg)\\
		&\stackrel{Ax1}{\equiv} (0\leq i <|s|\wedge 0\leq j+1\leq |s|)\wedge_L \Bigg(\Big((s[j+1]>s[i])\wedge (wp(i:=j,I_{j+1}^{j}))\Big) 
   				\\&\vee \Big(\neg  (s[j+1]>s[i])\wedge (I_{j+1}^{j})\Big)\Bigg)\\
		&\equiv   (0\leq i <|s|\wedge 0\leq j+1\leq |s|)\wedge_L 
		\\&\Bigg(\Big((s[j+1]>s[i])\wedge ((0\leq j <|s|\wedge 1\leq j+1\leq |s|)\wedge_L (\forall k:\mathbb{Z})(0\leq k<j+1\rightarrow_L s[k]\leq s[j])))\Big) 
   				\\&\vee \Big(\neg  (s[j+1]>s[i])\wedge ((0\leq i <|s|\wedge 1\leq j+1\leq |s|)\wedge_L (\forall k:\mathbb{Z})(0\leq k<j+1\rightarrow_L s[k]\leq s[i]))\Big)\Bigg)\\
	\end{align*}
	Una vez calculada la precondición más débil, debemos ver si $(I\wedge B)$ implican dicha precondición. Probamos cada
	parte por separado suponiendo que vale o no la guarda:
	\begin{itemize}
		\item $(0\leq i <|s|\wedge 0\leq j+1\leq |s|)\wedge_L $
		\\ $\Bigg(\Big((s[j+1]>s[i])\wedge ((0\leq j <|s|\wedge 1\leq j+1\leq |s|)\wedge_L (\forall k:\mathbb{Z})(0\leq k<j+1\rightarrow_L s[k]\leq s[j])))\Big)\Bigg)  $\smallskip \\
		Supongo que vale $(s[j+1]>s[i])$ y entonces vale lo restante por razones muy fuleras que no quiero demostrar .
		\item $ (0\leq i <|s|\wedge 0\leq j+1\leq |s|)\wedge_L$
		 \\ $\Bigg( \Big(\neg  (s[j+1]>s[i])\wedge ((0\leq i <|s|\wedge 1\leq j+1\leq |s|)\wedge_L (\forall k:\mathbb{Z})(0\leq k<j+1\rightarrow_L s[k]\leq s[i]))\Big)\Bigg)$\smallskip \\
		Supongo que vale $\neg (s[j+1]>s[i])$ y entonces vale lo restante por razones muy fuleras que no quiero demostrar.\smallskip \\
		Como $(I\wedge B) \Rightarrow wp(ciclo,I)$, podemos afirmar que el cuerpo del ciclo preserva el invariante.
	\end{itemize}
\end{enumerate}

\paragraph{Ejercicio 7\\}
\noindent
\textbf{proc sumaPosicionesImpares}(in $n: seq\langle\mathbb{Z}\rangle$, out $result: \mathbb{Z})\ \{$\smallskip \\
\hspace*{6mm}\textbf{Pre}$\{True \}$\\
\hspace*{6mm}\textbf{Post}$\{ result=\sum_{i=0}^{|s|-1}(\textrm{if}\ i\ mod\ 2=1\ \textrm{then}\ s[i]\ \textrm{else}\ 0\ \textrm{fi})\}$\\
$\}$\medskip\\

\noindent
\begin{enumerate}[label=\alph*)]
	\item 

\begin{lstlisting}
j:=0;
result:=0;

|\textbf{while}|(j<s.size()) |\textbf{do}|
  |\textbf{if}| (j mod 2 = 1) 
     result=:result +s[j]
  |\textbf{else}|
     skip
  |\textbf{endif}|;
   j:=j+1
|\textbf{endwhile}|
\end{lstlisting}
	\item
	
 Especificacion del ciclo:
	\begin{itemize}
		\item $P_c: result=0 \wedge j=0$
		\item $Q_c: result=\sum_{i=0}^{|s|-1}(\textrm{if}\ i\ mod\ 2=1\ \textrm{then}\ s[i]\ \textrm{else}\ 0\ \textrm{fi}))$
		\item $I\equiv 0\leq j\leq |s| \wedge result=\sum_{i=0}^{j-1}(\textrm{if}\ i\ mod\ 2=1\ \textrm{then}\ s[i]\ \textrm{else}\ 0\ \textrm{fi})$
		\item $f_v:n-i$
	\end{itemize}
	
	Solo demostraremos correción parcial de este ciclo.
	\begin{flushleft}$\scalebox{1.3}{$P_c\Rightarrow I$}$\end{flushleft}
	Debemos demostrar que vale $I$ sabiendo que vale $P_c$ 
	\begin{itemize}
		\item $ 0\leq j\leq |s|$\smallskip \\
		Por $P_c$ sabemos que $j=0$, entonces, la primera desigualdad es una tautologia,
		$0\leq 0$, para la segunda vemos que $|s|\geq 0$, luego como 
		$|s|$, no puede ser menor que cero es valido.
		\item  $result=\sum_{i=0}^{j-1}(\textrm{if}\ i\ mod\ 2=1\ \textrm{then}\ s[i]\ \textrm{else}\ 0\ \textrm{fi})$\\
		Trivial, por $P_c$ $result=0$, y  como $j=0$, entonces
		 $\sum_{i=0}^{-1}(\textrm{if}\ i\ mod\ 2=1\ \textrm{then}\ s[i]\ \textrm{else}\ 0\ \textrm{fi})=result=0$.
	\end{itemize}	 
	$\scalebox{1.3}{$(I\wedge \neg B)\Rightarrow Q_c$}$\smallskip \\
	Queremos demostrar que vale $Q_c$, asumiendo que vale  $I\wedge \neg B$.
	\begin{itemize}
	\item cuando deja de valer la guarda, $j=|s|$ sabemos por el invariante que $result=\sum_{i=0}^{j-1}(\textrm{if}\ i\ mod\ 2=1\ \textrm{then}\ s[i]\ \textrm{else}\ 0\ \textrm{fi})$, como $j=|s|$, entonces,
	vale que $result=\sum_{j=1}^{|s|-1}(\textrm{if}\ n\ mod\ j=0\ \textrm{then}\ j\ \textrm{else}\ 0\ \textrm{fi})$, que es exactamente lo mismo que $Q_c$
	\end{itemize}
	$\scalebox{1.3}{$\{I\wedge B\} \textbf{ ciclo } \{I\}$}$\medskip \\
	Queremos ver que vale la siguiente tripla de Hoare $\{I\wedge B\} \textbf{ ciclo } \{I\}$.\\
	Llamemos S1 a la primer instrucción del cuerpo del ciclo, S2 a la segunda:\\
	S1$: \textbf{if}\ (j\ mod\ 2=1\ \textbf{then}\ result:=result+s[j]\ \textbf{else}\ skip\ \textbf{endif};$\\
	S2$: j:=j+1$\\
	Lo primero que haremos es calcular $wp(ciclo,I)$.
	\begin{equation}wp(S1;S2,I)\stackrel{Ax3}{\equiv}wp(S1,wp(S2,I))\end{equation}
	Antes de seguir, debemos calcular $wp(S2,I)$. Para eso usaremos el axioma 1:
	\begin{align*}
	wp(S2,I)&\stackrel{Ax1}{\equiv}def(j+1)\wedge_L I_{j+1}^{j}\\
		&\equiv true\wedge_L 0\leq j+1\leq |s| \wedge result=\sum_{i=0}^{j-1+1}(\textrm{if}\ i\ mod\ 2=1\ \textrm{then}\ s[i]\ \textrm{else}\ 0\ \textrm{fi})\\
		&\equiv 0\leq j+1\leq |s| \wedge result=\sum_{i=0}^{j-1+1}(\textrm{if}\ i\ mod\ 2=1\ \textrm{then}\ s[i]\ \textrm{else}\ 0\ \textrm{fi}))
	\end{align*}
	
	Volviendo a (4), reemplazamos $wp(S2,I)$ y nos queda:
	\begin{align*}
	wp(S1;S2,I)&\stackrel{Ax3}{\equiv}wp(S1,wp(S2,I))\\
	&\equiv  wp(S1,0\leq j+1\leq |s| \wedge result=\sum_{i=0}^{j-1+1}(\textrm{if}\ i\ mod\ 2=1\ \textrm{then}\ s[i]\ \textrm{else}\ 0\ \textrm{fi}))\\
	&\equiv  wp(S1,I_{j+1}^{j})\\
	&\stackrel{Ax4}{\equiv}def(j\ mod\ 2=1)\wedge_L \Bigg(\Big((j\ mod\ 2=1)\wedge (wp(result:=result+s[j],I_{j+1}^{j}))\Big) 
   				\\&\vee \Big(\neg  (j\ mod\ 2=1)\wedge (wp(skip,I_{j+1}^{j}))\Big)\Bigg)\\
		&\stackrel{Ax1}{\equiv} True\wedge_L \Bigg(\Big((j\ mod\ 2=1)\wedge (wp(result:=result+s[j],I_{j+1}^{j}))\Big) 
   				\vee \Big(\neg  (j\ mod\ 2=1)\wedge I_{j+1}^{j})\Big)\Bigg)\\
		&\equiv  \Bigg(\Big((j\ mod\ 2=1)\wedge (wp(result:=result+s[j],I_{j+1}^{j}))\Big) 
   				\vee \Big(\neg  (j\ mod\ 2=1)\wedge I_{j+1}^{j})\Big)\Bigg)\\
   		&\equiv  \Bigg(\Big((j\ mod\ 2=1)\wedge (0\leq j+1\leq |s| \wedge result+s[j]=\sum_{i=0}^{j-1+1}
   			(\textrm{if}\ i\ mod\ 2=1\ \textrm{then}\ s[i]\ \textrm{else}\ 0\ \textrm{fi}))\Big) 
   				\\&\vee \Big(\neg  (j\ mod\ 2=1)\wedge 0\leq j+1\leq |s| \wedge result=\sum_{i=0}^{j-1+1}(\textrm{if}\ i\ mod\ 2=1\ \textrm{then}\ s[i]\ \textrm{else}\ 0\ \textrm{fi}))\Big)\Bigg)\\
	\end{align*}
	Una vez calculada la precondición más débil, debemos ver si $(I\wedge B)$ implican dicha precondición. Probamos cada
	parte por separado:
	\begin{itemize}
		\item $\Big((j\ mod\ 2=1)\wedge (0\leq j+1\leq |s| \wedge result+s[j]=\sum_{i=0}^{j-1+1}
   			(\textrm{if}\ i\ mod\ 2=1\ \textrm{then}\ s[i]\ \textrm{else}\ 0\ \textrm{fi}))\Big) $\smallskip \\
		Supongo que vale la izquierda y entonces vale la derecha.
		\item $\Big(\neg  (j\ mod\ 2=1)\wedge 0\leq j+1\leq |s| \wedge result=\sum_{i=0}^{j-1+1}(\textrm{if}\ i\ mod\ 2=1\ \textrm{then}\ s[i]\ \textrm{else}\ 0\ \textrm{fi}))\Big)$\smallskip \\
		Supongo que vale la izquierda y entonces vale la derecha.\smallskip \\
		Como $(I\wedge B) \Rightarrow wp(ciclo,I)$, podemos afirmar que el cuerpo del ciclo preserva el invariante.
	\end{itemize}
\end{enumerate}
%\includegraphics[width=\linewidth]{pe.png}


\end{document}